% SOURCE_AUTHORITY: SGA5 (1).pdf, PDF pages 3--4 (Roman pages III--IV; scan-only authority).
% SOURCE_SHA256: B256EBD072A8C68209518412A263C9289C6F1854A346733D86F885930D5FE6CA
\thispagestyle{empty}
\section*{\uline{Introdução}}

Este volume reúne as exposições do seminário de Grothendieck realizado no IHES em 1965/66, distribuídas inicialmente sob a forma de notas mimeografadas. Em relação à versão primitiva, as únicas mudanças importantes dizem respeito à exposição II, que não foi reproduzida, e à exposição III, que foi completamente reescrita e ampliada com um apêndice numerado III~B. Com exceção de algumas modificações de detalhe e do acréscimo de notas de rodapé, as demais exposições foram deixadas como estavam.

O núcleo do seminário é constituído pela fórmula de Lefschetz em cohomologia étale (III, III~B, XII) e por sua aplicação à interpretação cohomológica das funções $L$ (XIV). Todos os resultados anunciados por Grothendieck em sua exposição no seminário Bourbaki~[2] são aqui demonstrados por completo. As fórmulas de traço estabelecidas em (III, III~B e XII), por diferentes métodos, são mais gerais do que o necessário para demonstrar apenas a racionalidade das funções $L$. Na exposição de Deligne (SGA~$4^{1/2}$ Rapport) encontra-se uma demonstração completa desta última, consideravelmente mais curta, que segue o método de Grothendieck das exposições XII e XIV, mas sem toda generalidade supérflua. Esperamos, contudo, que as fórmulas de III e III~B possam servir em outras situações. O restante do seminário consiste em duas exposições sobre a teoria da passagem ao limite que dá origem à cohomologia $\ell$-ádica (V, VI) e em diversos complementos ao formalismo de dualidade (I, VII) e à fórmula de Lefschetz (VII, X).

Aqui está, com mais precisão, o conteúdo deste volume.

A exposição I é independente do restante do seminário. Seu principal resultado afirma que, sobre um esquema regular que satisfaz certas hipóteses adicionais de natureza local, e sempre que se disponha da resolução de singularidades e do teorema de pureza, o feixe constante $\mathbb Z/n\mathbb Z$, para $n$ primo com as características residuais, é dualizante (I.3.4.1). Esse teorema vale, em particular, para os esquemas regulares excelentes de característica zero, graças a Hironaka e aos resultados de Artin (SGA~4.XIX). O caso de um esquema regular de dimensão $1$ é

% SOURCE_PAGE_BOUNDARY: the printed sentence continues from Roman p. III to p. IV here.
\clearpage
\noindent tratado separadamente, por outro método, muito elementar (ver também (SGA~$4^{1/2}$ Dualité)). Observemos que, por meio de um argumento diferente, que não usa nem resolução nem pureza, Deligne demonstra em (SGA~$4^{1/2}$ Th. Finitude) que, se $f:X\longrightarrow S$ apresenta $X$ como um esquema de tipo finito sobre um esquema regular $S$ de dimensão $0$ ou $1$, o complexo $f^!(\mathbb Z/n\mathbb Z)_S$ (com $n$ primo com as características residuais) é dualizante. Não se sabe, contudo, se esse resultado, mais útil do que o teorema de bidualidade da exposição I, com o qual coincide parcialmente sem o conter, estende-se ao caso de um esquema-base excelente $S$, regular e de dimensão $>1$.

Como foi dito anteriormente, a exposição II, intitulada ``Fórmulas de Künneth para a cohomologia com suportes quaisquer'', não figura neste volume. Redigida por L.~Illusie a partir de notas manuscritas de Grothendieck, era dedicada a teoremas gerais sobre a propriedade de ser cohomologicamente próprio e sobre a aciclicidade local, mas estes só eram demonstrados sob a hipótese de resolução. Uma demonstração dos mesmos resultados sem essa hipótese, obtida posteriormente por Deligne (SGA~$4^{1/2}$ Th. Finitude), tornou desnecessária a publicação da exposição. Alguns complementos, relativos, por exemplo, ao caso dos divisores com cruzamentos normais, foram incorporados a um apêndice de (loc.~cit.). Acrescentemos que, na exposição da Sra.~Raynaud (SGA~1 XIII), demonstram-se, para feixes de conjuntos e de grupos não comutativos, teoremas análogos aos de (loc.~cit.).

A exposição III contém uma demonstração da fórmula de Lefschetz--Verdier para correspondências cohomológicas entre complexos de feixes sobre esquemas de tipo finito sobre um corpo; a exposição III~B mostra como calcular os termos locais da fórmula em certos casos, por exemplo para a correspondência de Frobenius sobre curvas. Remetemos o leitor às introduções das exposições III e III~B para mais detalhes sobre seu conteúdo.

A exposição IV, de A.~Grothendieck, dedicada à construção da classe de cohomologia associada a um ciclo, foi redigida por P.~Deligne e publicada em SGA~$4^{1/2}$.
